\documentclass[12pt,a4paper]{report}

\usepackage[utf8]{inputenc}
\usepackage[T1]{fontenc}
\usepackage[spanish,es-tabla]{babel}
\usepackage{../lab-doc}


\doctitulo{Manual Técnico del Portal del Laboratorio de Archivística y Humanidades Digitales}
\docsubtitulo{Manual técnico}
\docautores{Equipo del Laboratorio de Archivística\\ Universidad Distrital Francisco José de Caldas}
\docciudad{Bogotá, Colombia}
\docfecha{2026}

\begin{document}

\portadalab

\tableofcontents
\thispagestyle{chapternoline}

\chapter{Introducción y alcance}

El presente manual técnico documenta la arquitectura, el despliegue, la
gestión de contenido, el mantenimiento y la resolución de problemas del
\textbf{portal web del Laboratorio de Archivística y Humanidades Digitales} de
la Universidad Distrital Francisco José de Caldas.

El alcance de este documento se limita a los componentes que se encuentran
versionados en el repositorio \texttt{Lab-Archivistica-Frontend}:

\begin{itemize}
  \item Frontend desarrollado con \textbf{Astro} (sitio estático generado).
  \item Proxy \textbf{nginx} global que publica el portal en el puerto 80.
  \item Stack \textbf{OrfeoNG} (gestión documental).
  \item Stack \textbf{DSpace 10} (repositorio digital).
  \item Servicio \textbf{Apache Tika} (extracción de metadatos).
  \item Scripts de despliegue, operación y actualización de contenido.
  \item Manuales de usuario y técnico (directorio \texttt{docs/}).
\end{itemize}

Las aplicaciones \textbf{Archivematica}, \textbf{AtoM} y \textbf{ExifTool} se
despliegan fuera de este repositorio (en el servidor host) y se integran
mediante el proxy. Este manual las menciona a nivel de integración, pero no
documenta su instalación interna.

\section{Público objetivo}

Este manual está dirigido a:

\begin{itemize}
  \item \textbf{Administradores del portal}: responsables de desplegar,
        actualizar y mantener el sitio.
  \item \textbf{Gestores de contenido}: personas encargadas de publicar
        noticias y mantener los catálogos de equipos, salas y software.
  \item \textbf{Desarrolladores}: quienes modifican el código del frontend.
\end{itemize}

\section{Convenciones}

\begin{itemize}
  \item Los comandos de terminal aparecen en cajas de \textbf{Terminal}.
  \item Los fragmentos de código o configuración aparecen en cajas de
        \textbf{Código}.
  \item Las advertencias y recomendaciones aparecen en cajas resaltadas.
\end{itemize}

\chapter{Visión general de la arquitectura}

El portal sigue un modelo de \textbf{microservicios por contenedores}
orquestado con Docker Compose. Un único proxy \texttt{nginx}
(\texttt{lab-proxy}) actúa como puerta de entrada pública y distribuye el
tráfico hacia cada aplicación según la ruta solicitada.

\begin{center}
\codeblock[Diagrama de arquitectura]{
Usuario --> nginx (lab-proxy, puerto 80)\\
  /               -> astro-frontend       (4321)\\
  /orfeo/...       -> orfeo-nginx          (80)\\
  /dspace/         -> dspace-frontend      (4000)\\
  /dspace/server/  -> dspace-api           (8080)\\
  /tika/api/       -> apache-tika          (9998)\\
  /atom/           -> host.docker.internal:63001  (AtoM)\\
  /archivematica/  -> host.docker.internal:63002\\
  /exiftool/       -> host.docker.internal:63003\\
}
\end{center}

\section{Contenedores del despliegue}

\begin{center}
\small
\begin{tabular}{|l|l|l|}
\hline
\textbf{Contenedor} & \textbf{Imagen / Base} & \textbf{Puerto expuesto} \\
\hline
\texttt{lab-proxy} & nginx:1.25-alpine & 80 \\
\texttt{astro-frontend} & node:22-alpine + serve & 127.0.0.1:4321 \\
\texttt{orfeo-mysql} & mysql:5.7 & 127.0.0.1:3307 \\
\texttt{orfeo-backend} & build local & interno \\
\texttt{orfeo-frontend} & build local & interno \\
\texttt{orfeo-sockets} & build local & interno \\
\texttt{orfeo-nginx} & nginx:1.25-alpine & 127.0.0.1:64000 \\
\texttt{dspace-db} & dspace-postgres-loadsql:10 & interno \\
\texttt{dspace-solr} & dspace-solr:10 & interno \\
\texttt{dspace-api} & dspace:dspace-10\_x & 127.0.0.1:8080 \\
\texttt{dspace-frontend} & dspace-angular:10 & 127.0.0.1:4000 \\
\texttt{apache-tika} & apache/tika:latest-full & 127.0.0.1:9998 \\
\hline
\end{tabular}
\end{center}

\begin{mdframed}[style=importantstyle]
  \textbf{Nota:} los puertos 3307, 4321, 64000, 8080, 4000 y 9998 están
  vinculados únicamente a \texttt{127.0.0.1}; el único puerto expuesto al
  exterior es el \texttt{80} del proxy. Las aplicaciones Orfeo, DSpace y Tika
  se publican vía \texttt{lab-proxy}.
\end{mdframed}

\section{Red y comunicación entre contenedores}

Todos los contenedores del despliegue pertenecen a la red Docker
\texttt{lab-network} y se comunican entre sí por su nombre de servicio.
Las aplicaciones externas (AtoM, Archivematica, ExifTool) se alcanzan a
través de \texttt{host.docker.internal}, que resuelve la IP del host.

\begin{itemize}
  \item \texttt{astro-frontend} sirve el sitio estático generado por Astro
        (carpeta \texttt{dist}) usando \texttt{serve} en el puerto 4321.
  \item \texttt{lab-proxy} recibe el tráfico del puerto 80 y lo reenvía según
        la tabla de rutas del capítulo 6.
  \item \texttt{orfeo-nginx} sirve como proxy interno del stack OrfeoNG
        (backend PHP + frontend + sockets).
  \item \texttt{dspace-api} ejecuta la migración de base de datos al arrancar
        y expone la API REST de DSpace en el puerto 8080.
  \item \texttt{dspace-frontend} sirve la interfaz Angular 19 de DSpace en el
        puerto 4000.
\end{itemize}

\chapter{Stack tecnológico}

\section{Frontend (Astro)}

El portal es una aplicación \textbf{Astro} con componentes \textbf{React}:

\begin{itemize}
  \item Astro 7.1.6.
  \item React 19.
  \item Tailwind CSS v4 (paleta personalizada: \texttt{navy} y \texttt{accent}).
  \item TypeScript en modo estricto.
  \item \texttt{@tailwindcss/typography} para el contenido de las noticias.
  \item \textbf{Content collections} con el cargador \texttt{glob} para leer
        las noticias desde \texttt{public/news/}.
  \item \textbf{Astro transitions} (\texttt{ClientRouter}) para las
        transiciones entre páginas.
\end{itemize}

\subsection{Ejecución local del frontend}

Instalación de dependencias:

\singlecmd{npm install}

En modo desarrollo (servidor con recarga en caliente):

\singlecmd{npm run dev}

Para generar el sitio estático y previsualizarlo:

\multicmd{npm run build\\
npm run preview}

El sitio generado queda en \texttt{dist/} y es lo que sirve el contenedor
\texttt{astro-frontend}.

\section{Contenedor del frontend}

El \texttt{Dockerfile} usa \texttt{node:22-alpine} para construir el sitio y
sirve la carpeta \texttt{dist} con \texttt{serve} en el puerto 4321:

\codeblock[Etapas del Dockerfile]{
FROM node:22-alpine AS build\\
\# ... npm ci + npm run build ...\\

FROM node:22-alpine\\
\# ... npm install -g serve ...\\
CMD ["serve", "dist", "-l", "4321"]\\
}

\section{Stack OrfeoNG}

El stack OrfeoNG está compuesto por:

\begin{itemize}
  \item \texttt{orfeo-mysql}: base de datos MySQL 5.7 (volumen
        \texttt{orfeo\_mysql\_data}).
  \item \texttt{orfeo-backend}: API PHP del sistema.
  \item \texttt{orfeo-frontend}: interfaz de usuario.
  \item \texttt{orfeo-sockets}: servidor de WebSockets para notificaciones.
  \item \texttt{orfeo-nginx}: proxy interno del stack.
\end{itemize}

Los repositorios de OrfeoNG se clonan en tiempo de construcción desde GitHub
(variables \texttt{BACKEND\_REPO}, \texttt{FRONTEND\_REPO},
\texttt{SOCKETS\_REPO}).

\section{Stack DSpace 10}

\begin{itemize}
  \item \texttt{dspace-db}: PostgreSQL con carga inicial (SQL de DSpace 10).
  \item \texttt{dspace-solr}: índices Solr (búsqueda y descubrimiento).
  \item \texttt{dspace-api}: API REST de DSpace (Java, puerto 8080).
  \item \texttt{dspace-frontend}: interfaz Angular 19 (puerto 4000).
\end{itemize}

Los datos persistentes viven en los volúmenes \texttt{dspace\_pgdata},
\texttt{dspace\_solr\_data} y \texttt{dspace\_assetstore}.

\section{Apache Tika}

El contenedor \texttt{apache-tika} usa la imagen \texttt{apache/tika:latest-full}
y expone la API en el puerto 9998 (solo \texttt{127.0.0.1}). La interfaz web
estática se sirve desde \texttt{tika-ui/} a través del proxy.

\chapter{Estructura del repositorio}

\begin{center}
\small
\begin{tabular}{|l|l|}
\hline
\textbf{Ruta} & \textbf{Contenido} \\
\hline
\texttt{src/pages/} & Páginas Astro (Inicio, Equipos, Software, Salas, Noticias, Investigación) \\
\texttt{src/components/} & Componentes (Hero, Footer, SoftwareEquipment, etc.) \\
\texttt{src/layouts/} & Layout principal (encabezado, pie y metadatos) \\
\texttt{src/data/equipos.ts} & Datos centralizados de equipos \\
\texttt{src/content/} & Script \texttt{actualizar\_pagina-noticias.sh} \\
\texttt{src/content.config.ts} & Esquema de las content collections \\
\texttt{public/news/} & Noticias en Markdown (content collection) \\
\texttt{public/images/news/} & Imágenes de noticias \\
\texttt{public/} & Archivos estáticos (logos, favicon, robots.txt, sitemap) \\
\texttt{src/styles/global.css} & Estilos globales y paleta de colores \\
\texttt{nginx/default.conf} & Configuración del proxy \\
\texttt{docker-compose.yml} & Orquestación de todos los servicios \\
\texttt{orfeo/} & Código y configuración de OrfeoNG \\
\texttt{dspace/} & Configuración de DSpace (local.cfg, solr-init.sh) \\
\texttt{tika-ui/} & Interfaz estática de Tika \\
\texttt{Levantar\_Pagina.sh} & Rebuild del frontend \\
\texttt{gestor-app-lab.sh} & Levantar/derribar los 3 stacks \\
\texttt{docs/} & Manuales (usuario y técnico) en LaTeX \\
\hline
\end{tabular}
\end{center}

\section{Páginas del portal}

\begin{center}
\small
\begin{tabular}{|l|l|}
\hline
\textbf{Ruta} & \textbf{Contenido} \\
\hline
\texttt{/} & Página de inicio (carrusel, servicios, noticias, software, proceso, contacto) \\
\texttt{/equipos} & Catálogo de equipos con filtros \\
\texttt{/equipos/<id>} & Ficha de detalle de un equipo \\
\texttt{/software} & Catálogo de software con búsqueda y filtros \\
\texttt{/servicios} & Vista general de servicios \\
\texttt{/servicios/salas} & Listado de salas \\
\texttt{/servicios/salas/<id>} & Ficha de detalle de una sala \\
\texttt{/noticias} & Listado de noticias \\
\texttt{/noticias/<slug>} & Página de detalle de una noticia \\
\texttt{/investigacion} & Proyectos y líneas de investigación \\
\texttt{/404} & Página de error 404 \\
\hline
\end{tabular}
\end{center}

\chapter{Despliegue local}

\section{Requisitos}

\begin{itemize}
  \item Docker Engine 24+ y Docker Compose v2.
  \item Node.js 22 (solo para desarrollo del frontend).
  \item Git.
  \item Conexión a internet para la descarga de imágenes y repositorios de
        Orfeo (que se clonan en tiempo de construcción).
\end{itemize}

\section{Variables de entorno}

La configuración sensible se define mediante variables de entorno (archivo
\texttt{.env} en la raíz). Las principales son:

\begin{itemize}
  \item \texttt{MYSQL\_ROOT\_PASSWORD}, \texttt{DB\_NAME}, \texttt{DB\_USER},
        \texttt{DB\_PASS}: credenciales de MySQL de Orfeo.
  \item \texttt{AES\_KEY}: clave de cifrado compartida por Orfeo.
  \item \texttt{SERVER\_URL}: URL pública de Orfeo (por defecto
        \texttt{http://localhost:64000}).
  \item \texttt{PORT}: puerto host de \texttt{orfeo-nginx} (por defecto 64000).
  \item \texttt{BACKEND\_REPO}, \texttt{FRONTEND\_REPO},
        \texttt{SOCKETS\_REPO}: repositorios GitHub de OrfeoNG.
\end{itemize}

\begin{mdframed}[style=importantstyle]
  \textbf{Importante:} el archivo \texttt{.env} NO debe versionarse en el
  repositorio público. Copie \texttt{.env.example} (si existe) y ajuste los
  valores.
\end{mdframed}

\section{Levantar el portal completo}

\singlecmd{docker compose up -d}

Para construir las imágenes (primera vez o tras cambios en \texttt{Dockerfile}):

\singlecmd{docker compose build}

Para detener todos los servicios:

\singlecmd{docker compose down}

Para ver el estado de los contenedores:

\singlecmd{docker compose ps}

\section{Verificación del despliegue}

\begin{enumerate}
  \item Espere a que todos los contenedores estén en estado
        \texttt{running/healthy}: \singlecmd{docker compose ps}
  \item Verifique el portal principal:
        \singlecmd{curl -I http://localhost}
  \item Verifique las rutas de Orfeo, DSpace y Tika:
        \multicmd{curl -I http://localhost/orfeo/\\
        curl -I http://localhost/dspace/\\
        curl -I http://localhost/tika/}
  \item Revise los logs del proxy si alguna ruta falla:
        \singlecmd{docker compose logs -f lab-proxy}
\end{enumerate}

\chapter{Configuración del proxy nginx}

El archivo \texttt{nginx/default.conf} define las rutas del portal. A
continuación se resume la distribución del tráfico:

\begin{center}
\small
\begin{tabular}{|l|l|l|}
\hline
\textbf{Ruta} & \textbf{Destino} & \textbf{Observación} \\
\hline
\texttt{/} & astro-frontend:4321 & Portal principal \\
\texttt{/orfeo/} & orfeo-nginx:80 & UI de Orfeo \\
\texttt{/orfeo-ng/} & orfeo-nginx:80 & UI NG de Orfeo \\
\texttt{/dspace/} & dspace-frontend:4000 & UI Angular de DSpace \\
\texttt{/dspace/server/} & dspace:8080 & API REST de DSpace \\
\texttt{/server/} & dspace:8080 & Compatibilidad con instalaciones previas \\
\texttt{/ng\_backend/} & orfeo-nginx:80 & API de Orfeo \\
\texttt{/socket/}, \texttt{/socket.io/} & orfeo-nginx:80 & WebSockets de Orfeo \\
\texttt{/tika/} & tika-ui (estática) & Interfaz de Tika \\
\texttt{/tika/api/} & apache-tika:9998 & API de Tika \\
\texttt{/atom/} & host.docker.internal:63001 & AtoM (externo) \\
\texttt{/archivematica/} & host.docker.internal:63002 & Archivematica (externo) \\
\texttt{/exiftool/} & host.docker.internal:63003 & ExifTool (externo) \\
\hline
\end{tabular}
\end{center}

\begin{mdframed}[style=importantstyle]
  \textbf{Importante:} los bloques \texttt{/atom/}, \texttt{/archivematica/} y
  \texttt{/exiftool/} dependen de \texttt{host.docker.internal}, lo que exige
  que las aplicaciones correspondientes estén corriendo en el servidor host.
  Si alguna no está disponible, se generarán errores 502 en esas rutas.
\end{mdframed}

Para las rutas de Orfeo se usa un \texttt{resolver} de Docker
(\texttt{127.0.0.11}) con nombres de upstream variables; esto permite recargar
la configuración sin reiniciar el proxy cuando cambian las IPs de los
contenedores.

\section{Recargar la configuración del proxy}

Después de modificar \texttt{nginx/default.conf}, se recarga la configuración
sin detener el contenedor:

\singlecmd{docker compose exec lab-proxy nginx -s reload}

Para validar la sintaxis de la configuración antes de recargar:

\singlecmd{docker compose exec lab-proxy nginx -t}

\begin{mdframed}[style=importantstyle]
  \textbf{Recomendación:} pruebe siempre con \texttt{nginx -t} antes de
  recargar para evitar dejar el portal caído por un error de sintaxis.
\end{mdframed}

\chapter{Gestión de contenido}

La página se alimenta de distintas fuentes dentro del repositorio. Este
capítulo explica cómo agregar o modificar cada tipo de contenido: noticias,
equipos, salas, software e investigación. Salvo el script de noticias, los
cambios se aplican editando el código fuente y reconstruyendo el frontend con
\texttt{Levantar\_Pagina.sh} (Capítulo 8) o, en desarrollo, con
\texttt{npm run build}.

\section{Noticias}

Las noticias se escriben como archivos \textbf{Markdown} con portada
(\texttt{frontmatter}) y se almacenan en \texttt{public/news/}. El esquema se
define en \texttt{src/content.config.ts}:

\codeblock[content.config.ts]{
import \{ defineCollection, glob \} from 'astro:content';\\

const news = defineCollection(\{\\
  loader: glob(\{\\
    pattern: '**/[\textasciicircum\_]*.md',\\
    base: './public/news',\\
  \}),\\
  schema: z.object(\{\\
    title: z.string(),\\
    date: z.coerce.date(),\\
    excerpt: z.string(),\\
    author: z.string().default('GLUD'),\\
    image: z.object(\{ src: z.string(), alt: z.string() \}).optional(),\\
  \}),\\
\});\\

export const collections = \{ news \};\\
}

\subsection{Agregar una noticia}

\begin{enumerate}
  \item Crear un archivo \texttt{AAAA-MM-DD-slug.md} en \texttt{public/news/}.
        El nombre del archivo define la fecha de publicación y la URL de la
        noticia (\texttt{/noticias/AAAA-MM-DD-slug/}).
  \item Definir el \texttt{frontmatter}: \texttt{title}, \texttt{date},
        \texttt{excerpt}, \texttt{author} y, opcionalmente, \texttt{image}.
  \item Escribir el cuerpo en Markdown (se renderiza con
        \texttt{@tailwindcss/typography}).
  \item Ejecutar el script de actualización para regenerar la página.
\end{enumerate}

Ejemplo real de una noticia publicada:

\codeblock[public/news/2026-07-30-dspace10-disponible.md]{
---\\
title: "Nuevo DSpace 10 disponible y funcional"\\
date: 2026-07-30\\
excerpt: "El repositorio digital del laboratorio ha sido actualizado a DSpace 10."\\
author: "Nicolas Acevedo"\\
image: "/images/news/dspace.png"\\
---\\
}

Las imágenes se ubican en \texttt{public/images/news/} y se referencian con la
ruta absoluta \texttt{/images/news/<archivo>}.

\begin{mdframed}[style=importantstyle]
  \textbf{Nota:} la imagen es opcional. Si se omite, la tarjeta de la noticia
  no muestra imagen. Se recomienda PNG o SVG con fondo transparente para una
  mejor visualización.
\end{mdframed}

\subsection{Script de actualización}

El script \texttt{src/content/actualizar\_pagina-noticias.sh} reconstruye la
imagen \texttt{astro-frontend} y recrea el contenedor (sin reiniciar el resto
de servicios):

\singlecmd{bash src/content/actualizar\_pagina-noticias.sh}

Con reconstrucción limpia (sin caché):

\singlecmd{bash src/content/actualizar\_pagina-noticias.sh --no-cache}

\subsection{Orden de las noticias}

Las noticias se ordenan de forma descendente por su campo \texttt{date}. La
página de inicio muestra las 4 más recientes y el carrusel las 3 más
recientes. El listado completo (\texttt{/noticias}) las muestra todas.

\section{Equipos}

Los equipos se definen de forma centralizada en \texttt{src/data/equipos.ts}.
Este archivo alimenta la sección de la portada, el listado en
\texttt{/equipos} y la ficha de detalle en \texttt{/equipos/<id>}:

\codeblock[src/data/equipos.ts]{
export interface Equipo \{\\
  id: number;\\
  name: string;\\
  description: string;\\
  available: boolean;\\
  category: string;\\
\}\\

export const equipos: Equipo[] = [\\
  \{\\
    id: 1,\\
    name: 'Computadores',\\
    description: 'Equipo de cómputo de alta gama para procesamiento de datos y digitalización documental.',\\
    available: true,\\
    category: 'Cómputo',\\
  \},\\
];\\
}

\subsection{Agregar o modificar un equipo}

\begin{enumerate}
  \item Abra \texttt{src/data/equipos.ts}.
  \item Para agregar: añada un nuevo objeto al arreglo con un \texttt{id}
        consecutivo (los ids no deben repetirse, ya que definen la URL
        \texttt{/equipos/<id>}).
  \item Para modificar: edite los campos del objeto correspondiente.
  \item Para cambiar la disponibilidad: ajuste \texttt{available} a
        \texttt{true} o \texttt{false}.
  \item Reconstruya el frontend para publicar los cambios.
\end{enumerate}

\begin{mdframed}[style=importantstyle]
  \textbf{Nota:} a diferencia de las noticias, los cambios en equipos, salas y
  software requieren reconstruir el frontend para verse publicados.
\end{mdframed}

\section{Salas}

Las salas están definidas en el arreglo \texttt{salas} de
\texttt{src/pages/servicios/salas/[id].astro}:

\codeblock[{src/pages/servicios/salas/[id].astro}]{
const salas = [\\
  \{ id: 1, name: 'Sala de Digitalización', capacity: 6, available: true, description: '...' \},\\
  \{ id: 2, name: 'Sala de Trabajo Colaborativo', capacity: 12, available: true, description: '...' \},\\
  \{ id: 3, name: 'Cubículo de Estudio', capacity: 2, available: false, description: '...' \},\\
];\\
}

Cada sala define \texttt{id}, \texttt{name}, \texttt{capacity},
\texttt{available} y \texttt{description}. Al agregar o eliminar salas debe
actualizarse también \texttt{getStaticPaths()} con los \texttt{id}
disponibles:

\codeblock[getStaticPaths]{
export async function getStaticPaths() \{\\
  return [\\
    \{ params: \{ id: '1' \} \},\\
    \{ params: \{ id: '2' \} \},\\
    \{ params: \{ id: '3' \} \},\\
  ];\\
\}\\
}

\section{Software}

\subsection{Catálogo de software (\texttt{/software})}

El catálogo completo se define en el arreglo \texttt{softwares} de
\texttt{src/pages/software/index.astro}. Cada entrada define:

\begin{itemize}
  \item \texttt{id}: identificador numérico.
  \item \texttt{name}: nombre de la herramienta.
  \item \texttt{version}: versión mostrada.
  \item \texttt{category}: categoría (Repositorio Digital, Preservación
        Digital, Descripción Archivística, Gestión Documental, Extracción de
        Metadatos, Metadatos).
  \item \texttt{licenseType}: tipo de licencia.
  \item \texttt{description}: descripción breve.
  \item \texttt{url}: URL de acceso (opcional; si se omite, la tarjeta no
        muestra botón Acceder).
\end{itemize}

\subsection{Listado de software de la portada}

El listado de la portada se define en \texttt{src/components/SoftwareEquipment.astro}
en el arreglo \texttt{softwareList}. Cada entrada define \texttt{name},
\texttt{href} y \texttt{logo} (un SVG embebido):

\codeblock[src/components/SoftwareEquipment.astro]{
const softwareList = [\\
  \{ name: 'DSpace', href: '/dspace/', logo: '<svg ...>' \},\\
  \{ name: 'Archivematica', href: 'http://labhumanidades.udistrital.edu.co:62080', logo: '<svg ...>' \},\\
];\\
}

El \texttt{href} puede ser una ruta interna servida por el proxy nginx
(Capítulo 6) o una URL externa. El \texttt{logo} es un SVG embebido que se
renderiza con \texttt{set:html}.

\section{Investigación}

Los proyectos de investigación se definen en el arreglo embebido de
\texttt{src/pages/investigacion/index.astro}. Cada tarjeta define:

\begin{itemize}
  \item \texttt{title}: título del proyecto.
  \item \texttt{desc}: descripción breve.
  \item \texttt{tag}: etiqueta de estado (En curso, Investigación, Publicado).
  \item \texttt{color}: color de la etiqueta (green, blue, accent).
\end{itemize}

Para agregar un proyecto, añada un objeto al arreglo y reconstruya el frontend.

\section{Servicios de la portada}

Las tarjetas de servicios de la portada se definen en
\texttt{src/components/Services.astro} (Software especializado, Laboratorios
Virtuales, Préstamo de equipos, Préstamos de sala). Cada tarjeta define
\texttt{title}, \texttt{description}, \texttt{icon} (SVG) y \texttt{href}.

\section{Flujo de publicación de cambios}

\begin{enumerate}
  \item Edite el archivo correspondiente (Markdown de noticia, \texttt{equipos.ts},
        componente de salas/software/investigación).
  \item Pruebe en desarrollo: \singlecmd{npm run dev}
  \item Genere el sitio: \singlecmd{npm run build}
  \item Verifique en \texttt{dist/} que las páginas nuevas se generaron.
  \item En el servidor, actualice el repositorio con \texttt{git pull} y
        reconstruya el frontend.
\end{enumerate}

\chapter{Mantenimiento y operación}

\section{Levantar el frontend}

El script \texttt{Levantar\_Pagina.sh} reconstruye únicamente el contenedor
\texttt{astro-frontend} tras cambios en el frontend:

\singlecmd{bash Levantar\_Pagina.sh}

Para forzar una construcción limpia:

\singlecmd{bash Levantar\_Pagina.sh --no-cache}

\section{Gestor de los tres stacks}

El script \texttt{gestor-app-lab.sh} opera sobre los tres directorios del
servidor (\texttt{Lab-Archivistica-Frontend}, \texttt{archivematica/hack} y
\texttt{atom}):

\multicmd{bash gestor-app-lab.sh up\\
bash gestor-app-lab.sh down}

\section{Registros (logs)}

\codeblock[Registros por contenedor]{
docker compose logs -f astro-frontend\\
docker compose logs -f orfeo-backend\\
docker compose logs -f dspace\\
docker compose logs -f apache-tika\\
docker compose logs -f lab-proxy\\
}

Para ver los últimos N registros sin seguir:

\singlecmd{docker compose logs --tail 100 astro-frontend}

\section{Respaldos}

Los datos persistentes viven en volúmenes Docker:

\begin{itemize}
  \item \texttt{orfeo\_mysql\_data}: base de datos de Orfeo.
  \item \texttt{dspace\_pgdata}: base de datos de DSpace.
  \item \texttt{dspace\_solr\_data}: índices de Solr.
  \item \texttt{dspace\_assetstore}: objetos del repositorio.
\end{itemize}

\begin{mdframed}[style=importantstyle]
  \textbf{Recomendación:} realice respaldos periódicos de estos volúmenes
  (p.~ej. con \texttt{docker run --rm -v <volumen>:/data -v \$(pwd):/backup
  busybox cp -r /data /backup}) y guárdelos fuera del servidor.
\end{mdframed}

\subsection{Ejemplo de respaldo de un volumen}

\codeblock[Respaldo del volumen de Orfeo]{
docker run --rm -v orfeo\_mysql\_data:/data \\
  -v /ruta/respaldo:/backup busybox \\
  cp -r /data /backup/orfeo-mysql\\
}

\subsection{Actualización de imágenes base}

Para actualizar las imágenes base (nginx, node, DSpace, Tika):

\begin{enumerate}
  \item Haga un respaldo de los volúmenes de datos.
  \item Actualice las versiones en \texttt{docker-compose.yml}.
  \item Reconstruya: \singlecmd{docker compose build}
  \item Reinicie: \singlecmd{docker compose up -d}
  \item Verifique el estado y los datos después del reinicio.
\end{enumerate}

\section{Actualizar el código del repositorio}

\begin{enumerate}
  \item En el servidor, dentro de \texttt{Lab-Archivistica-Frontend}:
        \singlecmd{git pull}
  \item Si hubo cambios en \texttt{public/news/}: ejecute
        \singlecmd{bash src/content/actualizar\_pagina-noticias.sh}
  \item Si hubo cambios en el código del frontend (componentes, datos,
        estilos): \singlecmd{bash Levantar\_Pagina.sh}
  \item Si hubo cambios en \texttt{docker-compose.yml} o en el
        \texttt{Dockerfile}: \singlecmd{docker compose up -d --build}
  \item Si hubo cambios en \texttt{nginx/default.conf}: valide y recargue el
        proxy.
\end{enumerate}

\section{Monitoreo básico}

\begin{itemize}
  \item Estado de los contenedores: \singlecmd{docker compose ps}
  \item Uso de recursos: \singlecmd{docker stats}
  \item Espacio en disco: \singlecmd{df -h}
  \item Verificación de respuesta del portal: \singlecmd{curl -I http://localhost}
\end{itemize}

\chapter{Resolución de problemas}

\begin{itemize}
  \item \textbf{El puerto 80 está ocupado.} Detenga el servicio que lo usa o
        cambie el mapeo en \texttt{docker-compose.yml}.
  \item \textbf{Error 502 en \texttt{/atom/}, \texttt{/archivematica/} o
        \texttt{/exiftool/}.} Verifique que la aplicación externa esté
        corriendo en el host y que \texttt{host.docker.internal} resuelva.
  \item \textbf{Orfeo no responde.} Confirme que \texttt{orfeo-mysql} esté
        sano (\texttt{docker compose ps}) y revise los logs del backend.
  \item \textbf{DSpace tarda en iniciar.} La API ejecuta una migración de base
        de datos al arrancar; espere a que el healthcheck del contenedor pase.
  \item \textbf{Noticias desactualizadas.} Ejecute
        \texttt{actualizar\_pagina-noticias.sh} para reconstruir el frontend.
  \item \textbf{Los contenedores no resuelven entre sí.} Asegúrese de que todos
        pertenecen a la red \texttt{lab-network}.
  \item \textbf{El sitio muestra una versión antigua tras un cambio.}
        Reconstruya la imagen sin caché
        (\texttt{Levantar\_Pagina.sh --no-cache}) y reinicie el contenedor.
  \item \textbf{Una imagen de noticia no aparece.} Verifique que el archivo
        exista en \texttt{public/images/news/} y que la ruta en el
        \texttt{frontmatter} comience con \texttt{/images/news/}.
  \item \textbf{La página de un equipo devuelve 404.} El \texttt{id} en
        \texttt{getStaticPaths()} (o en el arreglo \texttt{equipos}) no
        coincide con la URL solicitada; reconstruya el sitio.
  \item \textbf{El proxy recibe 502 en todas las rutas.} Verifique que
        \texttt{lab-proxy} y los servicios de destino estén en la misma red y
        en ejecución.
\end{itemize}

\chapter{Seguridad y buenas prácticas}

\begin{itemize}
  \item No exponga los puertos internos (3307, 8080, 4000, 9998, 64000) a la
        red pública; mantenga el mapeo a \texttt{127.0.0.1}.
  \item Proteja las credenciales en el archivo \texttt{.env}; no lo versiones
        en el repositorio público.
  \item Actualice periódicamente las imágenes base (\texttt{nginx:1.25-alpine},
        \texttt{node:22-alpine}, DSpace, Tika) para corregir vulnerabilidades.
  \item Active HTTPS (certificado del servidor) para cifrar el tráfico del
        portal en producción.
  \item Haga respaldos automáticos de los volúmenes de datos y pruebe su
        restauración.
  \item Restrinja el acceso a los formularios de reserva y las cuentas de
        administrador de cada herramienta (DSpace, Orfeo, Archivematica, AtoM).
  \item Documente en este manual cualquier cambio de arquitectura relevante.
\end{itemize}

\section{Buenas prácticas de gestión de contenido}

\begin{itemize}
  \item Verifique siempre la ortografía y el formato en \texttt{npm run dev}
        antes de publicar una noticia.
  \item Use nombres de archivo \texttt{AAAA-MM-DD-slug.md} coherentes para
        las noticias.
  \item Mantenga los \texttt{id} de equipos y salas estables (no reutilice
        ids eliminados en el corto plazo para evitar enlaces rotos).
  \item Comprima y optimice las imágenes de noticias (PNG/SVG recomendado).
  \item Después de cada publicación, verifique el sitio en producción.
\end{itemize}

\end{document}